% paper/sections/06_scheme_per_variable.tex
%
% §6  方程式項別空間離散化 (per-term spatial discretization)
%      §6.0 (本ファイル冒頭) chapter intro + 変数別スキーム方針
%      §6.1 (06b_advection)       CLS ψ FCCD 面ジェット移流
%      §6.2 (06c_fccd_advection)  運動量 FCCD/UCCD6 移流
%      §6.3 (06d_viscous_3layer)  粘性 3 層
%      §6.4 (06d_viscous_3layer)  ρ/μ 物性値補間規則
%      §6.5 (forward ref to §8 PPE) 圧力 jump 取り扱い
%
% ═══════════════════════════════════════════════════════════════════════════════
\section{方程式項別空間離散化：CLS 移流・運動量移流・粘性項・物性値補間}
\label{ch:per_term_spatial}
% ═══════════════════════════════════════════════════════════════════════════════

\noindent\textbf{本章の位置付け：}
§~\ref{sec:CCD}（§4）で CCD/FCCD/UCCD6 の差分演算子を，
§~\ref{ch:cls_specific}（§5）で CLS framework のアルゴリズム
（$\psi\leftrightarrow\phi$ 対応・$\rho=\rho_g+(\rho_l-\rho_g)\psi$ 関係を含む）
を確立した．
本章では各方程式項
（CLS $\psi$ 移流・運動量 $u,v$ 移流・粘性 $\bnabla\cdot(2\mu\bm{D})$・
物性値 $\rho,\mu$・圧力/毛管 jump の面閉包）について演算子選択方針と
離散化詳細を扱う．続く §~\ref{ch:time_integration}（§7）で時間積分との結合を扱う．

% ─────────────────────────────────────────────────────────────────────────────
\subsection{方程式項別スキーム選択方針}
\label{sec:scheme_per_variable}
% ─────────────────────────────────────────────────────────────────────────────

本節の中心的主張は「差分演算子は\textbf{場の特性}に合わせて
選択せよ」である．
CCD の $\mathcal{O}(h^6)$ 精度はステンシル支持内で場が $C^\infty$ 滑らかである
仮定に依存するため，$C^0$ kink・$\delta$ 関数・step profile を持つ変数に
一律 CCD を適用すると\textbf{演算子の前提違反}が発散・膨張・スプリアス振動として
表出する．

\noindent\textbf{変数別割り当て}：
\begin{center}
\begin{tabularx}{\linewidth}{@{}>{\raggedright\arraybackslash}p{0.18\linewidth}>{\raggedright\arraybackslash}p{0.18\linewidth}>{\raggedright\arraybackslash}X@{}}
\hline
変数 & 特性 & 推奨スキーム \\
\hline
$\psi$（CLS）& $C^0$ 急 tanh & FCCD 面ジェット移流（面値保存型；\S\ref{sec:advection_motivation}） \\
$u, v$ bulk & $C^\infty$ 滑 & \textbf{UCCD6}（\S\ref{sec:uccd6_def}；FCCD 面フラックス形は \S\ref{sec:fccd_advection} の面評価位置拡張） \\
$u, v$ 界面帯 & $\bnabla u$ kink & 相内 2 次中心 / 片側応力評価（\S\ref{sec:viscous_3layer}） \\
$\phi$（再初期化）& Eikonal & \textbf{Ridge--Eikonal}（\S\ref{sec:xi_sdf_reinit}）／ Godunov・WENO--HJ~\cite{JiangPeng2000}（基底 Eikonal solver） \\
$p$ / 毛管 jump & 界面で jump & Balanced--Force 面評価位置（\S\ref{sec:balanced_force}）+ affine-jump PPE / 変分毛管 Hodge 閉包（\S\ref{sec:split_ppe_framework}）\\
$\rho, \mu$ & step profile & 低次（$\rho,\mu$ は $\psi$ から代数更新；face/corner 平均は拡張規則；\S\ref{sec:rho_mu_interpolation}）\\
\hline
\end{tabularx}
\end{center}

\noindent\textbf{設計原則}：
\begin{itemize}
  \item \textbf{$\psi$ に節点 CCD を適用してはならない．}
        tanh プロファイルを節点 CCD で微分すると急勾配帯の曲率が増幅され，
        半陰的処理でも不安定化することを本稿の数値実験で確認した．
        本稿では FCCD 面値演算子（\S\ref{sec:fccd_def}）で面フラックスを評価することで
        振動源を切り，FCCD 面ジェット移流を標準経路とする．
        DCCD/節点中心比較経路の散逸設計は \S~\ref{sec:dccd_derivation}（§4.4）に集約．
  \item \textbf{$p$ と毛管 jump は節点勾配では閉じない．}
        圧力跳躍を単一節点場として滑らかに内挿すると，
        実効的に全域 $\mathcal{O}(1)$ のスプリアス圧力勾配を生む．
        本稿では §~\ref{sec:balanced_force} の同一面評価位置原理に従い，
        一括 CSF 縮約経路では圧力勾配と毛管体積力を同じ面演算子で評価し，
        分相経路では §~\ref{sec:split_ppe_framework} の affine-jump 面共鎖と
        表面エネルギーの仮想仕事から得る，圧力仕事と随伴な Hodge 閉包で扱う．
  \item \textbf{$\rho, \mu$ に CCD を適用してはならない．}
        step プロファイルを CCD で差分すると Gibbs 振動が生じ，
        物性値が $[\rho_\text{min},\rho_\text{max}]$ 範囲外になる．
        物性値は常に低次で評価する：$\rho^{n+1} = \rho_g + (\rho_l-\rho_g)\psi^{n+1}$
        および $\mu^{n+1} = \mu_g + (\mu_l-\mu_g)\psi^{n+1}$ の代数更新を基準とする．
        face/corner 平均は粘性フラックス拡張で必要になる場合の補助規則である．
        具体的な補間式は \S~\ref{sec:rho_mu_interpolation}（§6.4）に集約する．
  \item \textbf{bulk 速度の標準移流は UCCD6．}
        $u, v$ は $C^\infty$ 滑らか（界面帯を除く）．
        本稿では移流の風上安定化を演算子内部で担う
        UCCD6（\S\ref{sec:uccd6_def}）を標準経路とする．
        FCCD 面フラックス形は，縮約 CSF 経路では圧力・毛管・対流フラックスを揃え，
        標準 pressure-jump 経路では毛管面共鎖と対流面フラックスの評価位置を比較する
        面閉包として \S\ref{sec:fccd_advection} に定義する．
  \item \textbf{界面帯の速度勾配は相内応力評価で閉じる．}
        $\bnabla u$ の kink を CCD stencil が跨ぐと $\mathcal{O}(1)$ 振動．
        $|\phi|\le 3h$ または phase-indicator mismatch で帯を同定し，
        法線方向は相内 2 次中心 / 片側，接線方向は滑らかさが保たれる場合だけ
        CCD を使う（\S\ref{sec:viscous_3layer}）．
\end{itemize}

\noindent\textbf{本章の構成}：
\S~\ref{sec:advection_motivation}（§6.1）で CLS $\psi$ 用 FCCD 面ジェット移流，
\S~\ref{sec:fccd_advection}（§6.2）で運動量 FCCD/UCCD6 移流，
\S~\ref{sec:viscous_3layer}（§6.3）で粘性 3 層アーキテクチャ（Layer A/B/C），
\S~\ref{sec:rho_mu_interpolation}（§6.4）で $\rho,\mu$ 物性値補間規則を扱う．
圧力 $p$ と毛管 jump の閉包は，Balanced--Force の同一面評価位置原理を
\S~\ref{sec:balanced_force}（§8）で述べ，
分相 affine-jump PPE と変分毛管 Hodge 閉包を
\S~\ref{sec:split_ppe_framework}（§9）で定式化する．
時間積分との結合（演算子分裂順序・粘性陰解法・物理 $\Delta t$ CFL）は
\S~\ref{ch:time_integration}（§7）で扱う．

% ═══════════════════════════════════════════════════════════════════════════════
